\section{Introduction}
The \textit{LiveTree} project involves the operation of a demonstrator that measures the real-time energy consumption and the photovoltaic production of three buildings. To manage these resources, an Artificial Intelligence (AI) neural network model is utilized to predict the energy consumption for each building. These predictions are generated daily at midnight for the following day, using the consumption data from the last seven days as input.

\subsection{Problem statement}
The project faces a fundamental challenge concerning the \textbf{continuity of the data acquisition chain}, which is prone to measurement failures or sensor outages that can last anywhere from minutes to several hours. While the existing software architecture is designed to handle brief interruptions, it currently \textbf{does not guarantee continuous data supply} over the critical 7-day input period required by the AI prediction model.\\

The core issue is that the AI model is constrained and \textbf{cannot generate forecasts if there is any missing data within this 7-day selection}. This constraint severely \textbf{limits the overall reliability and performance} of the prediction system. Consequently, a primary objective of the research is the \textbf{ability to reconstitute this missing data} when sensors fail to acquire real data.\\

The problem cannot be solved using \textbf{naive} methods, such as simply averaging past data. This is because energy consumption and production exhibit a \textbf{strong seasonal and cyclical component}. For instance, consumption can drastically decrease by a factor of 100 during holidays. Therefore, the solution requires finding \textbf{specific algorithms} that can incorporate historical data, seasonal trends, and current parameters into the data generation process.

\subsection{Project objectives and proposed approach}
To mitigate this critical limitation and ensure \textbf{system resilience}, the project proposes\textbf{ adding a pilot module} to the existing architecture. The primary role of this new layer is to ensure the \textbf{continuity} of the \textbf{input data} by maintaining a backup of data that is correctly correlated with the season and the calendar.\\

The pilot module is designed to achieve the following objectives:
\begin{itemize}
    \item \textbf{Assure data continuity} by generating and injecting backup data in real time when measurements are missing.
    \item \textbf{Maintain reliable forecasts} by ensuring the forecasting model always receives a complete seven-day input window.
    \item \textbf{Reinforce system resilience} by preserving the coherence of reconstructed data with seasonality and calendar constraints.
\end{itemize}

\subsection{Utility for end users}
For end users, the pilot module improves service continuity: forecasts remain available even when sensors temporarily fail. This increases operational robustness, reduces manual intervention, and enables more consistent downstream use of predictions (monitoring dashboards, alerts, and energy optimization workflows).




%%%%%%%% TODO : PLACEHOLDER, TO MODIFY ONCE THE STRUCTURE IS ACTUALLY COMPLETED
\subsection{Report structure}
The remainder of this report is organized as follows: Section~\ref{sec:background} summarizes the technical background and constraints of the demonstrator and forecasting pipeline; Section~\ref{sec:relatedwork} reviews relevant state-of-the-art methods for handling missing time-series data in forecasting contexts; Section~\ref{sec:method} details the proposed pilot module and the reconstruction strategies considered; Section~\ref{sec:evaluation} presents evaluation protocols and metrics to quantify both reconstruction quality and forecast impact; finally, Section~\ref{sec:conclusion} concludes and outlines perspectives for deployment and future improvements.
%%%%%%%% END PLACEHOLDER




\section{State of the Art}
\subsection{Context and operational requirements}
LiveTree operates an energy demonstrator monitoring real-time electricity consumption and photovoltaic production for multiple buildings. Forecasts are produced daily at midnight using a fixed historical input window (the last seven days). In this operational setting, sensor or communication failures can create gaps ranging from minutes to hours; however, the forecasting model is constrained and cannot run if any point in the required input window is missing, which directly threatens service continuity and system reliability. The project therefore motivates the addition of a \emph{pilot/backup module} that detects outages and injects reconstructed values in real time to preserve a complete input window and maintain uninterrupted forecast production.%
% 

%%%%%%% TODO: MAYBE REWRITE
This framing shifts the problem from "best-effort imputation" to \emph{production-grade reconstruction}: methods must (i) handle long contiguous gaps, (ii) preserve intraday structure relevant for forecasting, (iii) incorporate strong calendar/seasonal effects (e.g., closures and holidays), and (iv) remain auditable and computationally bounded in a continuous pipeline.%
% 
%%%%%%%%%%%%%%%%%%%%%%%%

\subsection{Missing-data mechanisms in energy monitoring}
Energy monitoring outages are rarely Missing Completely At Random (MCAR). Missingness often correlates with operational regimes (maintenance windows, nights/weekends, network conditions), producing \emph{structured} or \emph{block} missingness. This matters because many reconstruction methods assume local continuity or randomly scattered gaps and degrade sharply when missingness is long, clustered, or periodic. Classical missing-data theory emphasizes that naive approaches can introduce systematic bias when the missingness mechanism is not random.%
% 

In addition, energy signals exhibit strong multi-scale seasonality (daily/weekly/annual) and abrupt regime changes (holidays, closure days, unusual events). Reconstructions that over-smooth peaks or ignore context can be numerically plausible yet operationally misleading, affecting the downstream behavior of the prediction model.%
% 

\subsection{Baseline deterministic approaches}
\subsubsection{Direct interpolation and smoothing}
Linear/spline interpolation and local smoothing are widely used baselines for short gaps due to simplicity and negligible compute. However, interpolation error increases rapidly with gap length, and performance can collapse for multi-hour/multi-day outages because these methods rely on local continuity. They also tend to attenuate ramps and peaks (intraday variability), which are often critical features for short-horizon energy forecasts.%
% 

\subsubsection{Seasonal decomposition + component-wise interpolation}
A more structured baseline decomposes the series into trend, seasonal, and residual components (e.g. STL/LOESS), imputes missing values at the component level, and recombines them. This explicitly preserves dominant periodic shape and typically outperforms naive interpolation for \emph{short} missingness.\\%
% 

Nevertheless, decomposition-based methods inherit important limitations: STL requires sufficient surrounding valid data to estimate components; long contiguous gaps can make decomposition unstable or infeasible, so can the lack of post-gap data (which is likely to occur in practical use-cases). Moreover, component-wise interpolation often yields over-smoothed reconstructions that preserve average daily shape but suppress sharp transitions, and it does not naturally enforce physical/operational constraints (non-negativity, caps, closure/holiday regimes) without additional rules. Consequently, recent practice positions decomposition as a helpful \emph{building block} rather than a standalone solution in operational pipelines.%
% 

\subsection{State-space and structural time-series models (Kalman family)}
\subsubsection{Principle and relevance to imputation}
State-space models represent observations $y_t$ via latent states $\alpha_t$:
$$
\alpha_t = F_t \alpha_{t-1} + \eta_t,\quad y_t = H_t \alpha_t + \epsilon_t,
$$
with Gaussian noises $\eta_t,\epsilon_t$ in the classical formulation. The Kalman filter/smoother provides optimal (linear-Gaussian) recursive inference, and missing observations can be handled naturally by skipping measurement updates (filtering) and using smoothing for retrospective reconstruction. Foundational work established the Kalman filter and its use for filtering/prediction in linear systems.%

Structural time-series models (also called unobserved-components models) further decompose $y_t$ into interpretable latent components (level/trend/seasonal/regression), making them attractive for energy signals where interpretability and modular inclusion of temperature and calendar effects are important. Comprehensive treatments of such models and their estimation via Kalman methods are given by Harvey and by Durbin--Koopman.%

Parameter estimation for state-space models with missing observations is often performed via maximum likelihood using the EM algorithm combined with Kalman smoothing, providing a principled approach when the dynamics are reasonably modeled as linear and stochastic.%

\subsubsection{Strengths}
\begin{itemize}
\item \textbf{Probabilistic reconstruction:} yields uncertainty estimates, enabling "confidence-aware" backup injection and anomaly flagging.
\item \textbf{Context integration:} regression components can incorporate exogenous variables (e.g. temperature) and calendar indicators (holiday/closure).
\item \textbf{Online operation:} filtering is sequential and can be run continuously with bounded per-step cost \emph{given a fixed state dimension}.
\end{itemize}

\subsubsection{Limitations in high-frequency, multi-seasonal settings}
The key practical bottleneck is that classical seasonal components can require a large state dimension proportional to the seasonal period (e.g., dummy-seasonal models with $S{-}1$ states). For high-frequency series (e.g., 10-minute data) and multiple seasonalities (daily + weekly), naive structural formulations can become memory- and compute-intensive, which is consistent with implementation difficulties observed in practice for large seasonality settings.\\

A major line of work mitigates this by using \emph{trigonometric/Fourier seasonal representations} and innovations state-space formulations (e.g., TBATS) to represent long seasonal periods with far fewer states, enabling complex seasonality without quadratically progressing RAM usage.\\%

More generally, state-space methods can underperform when dynamics are strongly nonlinear, when regimes change abruptly (e.g., exceptional closures), or when missingness is long and coincides with regime transitions (the model may extrapolate the wrong regime). These limitations motivate hybrid approaches combining model-based smoothing with context-aware templates and constraints.

\subsection{Similarity-based methods and "analogue day" templates}
Similarity-based approaches reconstruct missing segments by borrowing patterns from historical periods with comparable context (same weekday/season, similar temperature regime, similar operational status) and aligning/scaling them to match the local level around the gap. This family includes $k$NN/donor-based imputation and pattern matching at daily or weekly granularity.\\

Within LiveTree-related work, WNNmulti (Weighted Nearest Neighbours in multivariate form) is described as a forecasting method that searches historical patterns across multiple related series (buildings and photovoltaic production), then produces predictions as a weighted average of the neighbors' future trajectories. It is computationally efficient and can exploit multivariate dependencies, but it does not solve the continuity problem by itself because it assumes a complete input window and does not reconstruct missing points \emph{inside} that window.\\%
%

For imputation, the same "nearest historical pattern" idea is often adapted as an \emph{analogue-day donor} mechanism: select candidate days/weeks under similar calendar and weather conditions, then reconstruct the missing interval using time-slot templates with scaling and optional blending at edges. Compared to purely statistical decomposition, donor/template methods remain usable even when decomposition fails (e.g., long gaps) and can preserve realistic intraday shapes if the donor selection is context-accurate. However, they can fail under rare events (no good historical analogue) and require robust similarity metrics and guardrails to prevent implausible injections.

\subsection{Hybrid and multi-layer reconstruction pipelines}
Recent operational consensus is that no single method dominates across all gap types; effective systems combine complementary strategies. In particular, decomposition is valuable for short gaps, while template/donor strategies and cross-series borrowing are more robust for long outages; both benefit from constraint enforcement and explicit logging.\\

A LiveTree-oriented "multi-layered" reconstruction logic is described as explicitly addressing decomposition weaknesses by combining multiple methods, adapting to gap length, applying contextual scaling/guardrails, and leveraging cross-site correlations, while logging donor choices and transformations for auditability.%
%

\begin{figure}[h]
\centering
\fbox{\begin{minipage}{0.95\linewidth}
\small
\textbf{Typical hybrid reconstruction flow:};
Gap detection $\rightarrow$ gap characterization (length, regime) $\rightarrow$
candidate generation (decomposition / donor templates / state-space smoothing / cross-series) $\rightarrow$
local scaling and constraint checks $\rightarrow$
injection + logs (donors, scaling, confidence).
\end{minipage}}
\caption{Operational view of hybrid imputation for continuous forecasting.}
\end{figure}

\subsection{Machine-learning imputation methods}
Beyond classical statistics, modern literature includes supervised and deep-learning imputers that learn nonlinear dependencies across variables and time. RNN-based methods (e.g., GRU-D, BRITS) incorporate missingness masks and time-since-last-observation to exploit informative missingness and multivariate structure.\\%
Attention-based architectures (e.g., SAITS) improve long-range dependency modeling and feature correlation handling.\\%
Probabilistic diffusion-based imputers (e.g., CSDI) further provide distributional reconstructions conditioned on observed context.\\%

While these methods can achieve strong benchmark accuracy, they introduce operational challenges for LiveTree-like constraints: substantial training and monitoring effort, lack of accounting for drift across seasons/years, harder explainability, and potentially high inference cost. They are therefore often integrated as optional components rather than the sole online solution, consistent with survey findings that practical systems are hybrid.%
% 

\subsection{Evaluation methodologies for imputation in continuous forecasting}
Imputation quality must be evaluated along two axes:
\begin{enumerate}
\item \textbf{Reconstruction fidelity:} MAE/RMSE (and scaled measures such as MASE) computed by artificially masking observed data under realistic missingness patterns, then comparing imputations to ground truth.
\item \textbf{Downstream impact:} the forecasting error induced by imputation, evaluated via backtesting in which forecasts are produced using partially reconstructed input windows.
\end{enumerate}
Forecast-accuracy metrics and rolling evaluation are well established in the forecasting literature, including scaled measures designed for comparability.%\\

In operational settings, evaluation should stratify by gap length, regime (weekday/weekend, holiday/closure), season, and "distance" to nearest real observations (end-of-window effects), because the hardest cases often occur near the most recent days where no future observations exist to anchor scaling.

\subsection{Synthesis: strengths, limitations, and gaps}
Table~\ref{tab:compare} summarizes how major approach families align with LiveTree constraints.\\




The main gaps in litterature regarding data reconstruction are as follows: (i) robust handling of long contiguous gaps under multi-scale seasonality without excessive memory, (ii) context-aware donor selection and scaling that remains stable at the end of the input window, (iii) principled use of cross-series information (other buildings / Solar panels) during outages, and (iv) auditable decision logs and safety constraints suitable for deployment. This further incentivizes hybrid pilot module combining gap-aware strategy selection, efficient seasonality representations (e.g., Fourier/TBATS-style), and constraint-driven scaling with explicit traceability.%
%
